\documentclass[11pt,a4paper]{article}
\usepackage[spanish]{babel}
\usepackage[utf8]{inputenc}
\usepackage[T1]{fontenc}
\usepackage{geometry}
\geometry{margin=2.4cm}
\usepackage{booktabs}
\usepackage{array}
\usepackage{amsmath}
\usepackage{xcolor}
\usepackage{hyperref}
\hypersetup{colorlinks=true, linkcolor=black, urlcolor=blue!50!black}
\usepackage{fancyhdr}
\pagestyle{fancy}
\fancyhf{}
\fancyhead[L]{\small GCI Credit Default Risk 2026}
\fancyhead[R]{\small Modelos del proyecto}
\fancyfoot[C]{\thepage}
\setlength{\parskip}{4pt}

\title{\textbf{Modelos del proyecto}\\[4pt]
\large GCI Global 2026 --- Home Credit Default Risk (subset de una tabla)}
\author{Mario Alejandro Marroqu\'in Mel\'endez}
\date{Octubre 2026}

\begin{document}
\maketitle

\begin{abstract}
\noindent Documento t\'ecnico de referencia sobre los modelos a utilizar en la
competencia GCI Global 2026 (Home Credit Default Risk, versi\'on de una sola
tabla). Se describe el baseline vigente, los modelos candidatos del ensemble,
la arquitectura de combinaci\'on y el protocolo de validaci\'on y selecci\'on.
La m\'etrica oficial es ROC-AUC; cierra el 20 de noviembre de 2026 (01:00 UTC)
y \textbf{solo cuenta la \'ultima submission} para el ranking final.
\end{abstract}

\tableofcontents

\section{Contexto y datos}
La competencia usa un \'unico archivo a nivel solicitud (\texttt{train.csv},
171{,}202 $\times$ 34; \texttt{test.csv}, 61{,}500 $\times$ 33), sin las tablas
auxiliares del Kaggle original. El objetivo \texttt{TARGET} est\'a desbalanceado:
8.1\% positivos. Hallazgos relevantes del EDA inicial:
\begin{itemize}
  \item \texttt{EXT\_SOURCE\_1/2/3} concentran la se\~nal (correlaci\'on de
        $-0.154$ a $-0.178$ con \texttt{TARGET}); sus valores faltantes predicen
        (default 8.9\% con \texttt{EXT\_SOURCE\_3} ausente vs.\ 7.7\% presente)
        $\Rightarrow$ no imputar; conservar NaN y agregar flags \texttt{*\_ISNA}.
  \item \texttt{DAYS\_EMPLOYED = 365243} es un c\'odigo especial presente en el
        18.0\% de las filas $\Rightarrow$ flag binario y conversi\'on a NaN.
  \item \texttt{ORGANIZATION\_TYPE} y \texttt{OCCUPATION\_TYPE} tienen alta
        cardinalidad $\Rightarrow$ target encoding \emph{out-of-fold} con
        smoothing (nunca global, por riesgo de leakage).
\end{itemize}

\section{Modelos}

\subsection{LightGBM --- modelo base (en producci\'on)}
Baseline actual del repositorio (commit del 2026-10-08): validaci\'on cruzada
estratificada de 5 folds, \textbf{CV AUC = 0.7581}.

\begin{table}[h]
\centering
\begin{tabular}{ll}
\toprule
\textbf{Hiperpar\'ametro} & \textbf{Valor} \\
\midrule
\texttt{learning\_rate} & 0.03 \\
\texttt{num\_leaves} & 31 \\
\texttt{min\_child\_samples} & 50 \\
\texttt{feature\_fraction} & 0.8 \\
\texttt{bagging\_fraction} / \texttt{bagging\_freq} & 0.8 / 1 \\
\texttt{lambda\_l2} & 1.0 \\
Early stopping & 200 rondas \\
Semillas & fijas en \texttt{src/config.py} \\
\bottomrule
\end{tabular}
\end{table}

\noindent\textbf{Rol:} baseline de referencia y candidato a ganador individual.
Ventajas: r\'apido en CPU, maneja NaN de forma nativa (clave para
\texttt{EXT\_SOURCE}), ideal para iterar feature engineering.

\subsection{XGBoost --- diversificador 1}
Mismo conjunto de features que LightGBM, pero con esquema de regularizaci\'on
distinto, lo que produce errores descorrelacionados --- condici\'on necesaria
para que el blend supere a los individuales. \textbf{Rol:} segundo pilar del
ensemble. Branch: \texttt{model/xgboost}.

\subsection{CatBoost --- diversificador 2}
Maneja variables categ\'oricas de forma nativa mediante \emph{ordered target
statistics} (anti-leakage incorporado), complementario al target encoding OOF
manual. Fuerte en datos tabulares sin ajuste fino. \textbf{Rol:} tercer pilar;
captura la informaci\'on de las categ\'oricas crudas. Branch:
\texttt{model/catboost}.

\subsection{TabPFN --- experimental (M4)}
\emph{Foundation model} tabular (Nature, 2025): aprende de manera
cualitativamente distinta a los \'arboles, lo que aporta diversidad real al
ensemble. Se eval\'ua en el Sprint 4; \textbf{solo se adopta si corre en CPU y
suma $\geq 0.002$ de AUC (RSKF) al blend}.

\section{Arquitectura del ensemble}
\begin{enumerate}
  \item \textbf{Nivel 0:} cada modelo se entrena con StratifiedKFold de 5 folds
        y produce predicciones \emph{out-of-fold} (OOF) sobre train y
        predicciones promediadas sobre test.
  \item \textbf{Nivel 1 --- blend ponderado:} combinaci\'on lineal de las OOF;
        los pesos se optimizan sobre las OOF (rejilla $0$--$1$, paso $0.1$).
  \item \textbf{Nivel 2 --- stacking (Sprint 4):} regresi\'on log\'istica sobre
        las OOF como meta-features.
  \item \textbf{Salida:} probabilidad final por fila $\rightarrow$ CSV
        \texttt{SK\_ID\_CURR,TARGET} en el formato exacto de
        \texttt{sample\_submission.csv}.
\end{enumerate}
Regla de adopci\'on: el blend/stacking solo reemplaza al mejor modelo individual
si lo supera en CV --- la complejidad adicional tambi\'en cuesta
reproducibilidad.

\section{Validaci\'on y selecci\'on}
\begin{itemize}
  \item \textbf{Desarrollo diario:} StratifiedKFold 5-fold con semilla fija;
        m\'etrica de control: AUC sobre OOF.
  \item \textbf{Selecci\'on del modelo final:} Repeated StratifiedKFold $3
        \times 5$ (semillas 42/43/44), decisi\'on por media $\pm$ desviaci\'on.
        \textbf{Nunca} por el public leaderboard (es un subconjunto del test y
        lleva a sobreajuste).
  \item \textbf{Se\~nal de alarma:} si $|\text{CV} - \text{LB}| > 0.02$, se
        detiene la iteraci\'on y se corrige primero la validaci\'on (domain
        shift o leakage). Complemento en M4: \emph{adversarial validation}
        train-vs-test, AUC $< 0.6$ = aceptable.
  \item \textbf{Umbrales de adopci\'on a \texttt{main}:} $+0.003$ features,
        $+0.002$ tuning, $+0.001$ blend/stacking. Un branch por experimento;
        \texttt{main} = mejor verificado; todo modelo guarda OOF, par\'ametros
        y semilla.
\end{itemize}

\section{Matem\'aticas del proyecto (para auditor\'ia)}

Esta secci\'on fija la formulaci\'on exacta de cada componente, de modo que
cualquier revisor pueda auditar el pipeline sin ambig\"uedades. Notaci\'on:
$N$ filas de train, $y_i \in \{0,1\}$, $x_i \in \mathbb{R}^d$, predicci\'on
$\hat{p}_i = P(y_i=1 \mid x_i)$.

\subsection{M\'etrica oficial: ROC-AUC}
\begin{equation}
\mathrm{AUC} = P\big(\hat{p}_i > \hat{p}_j \mid y_i=1,\ y_j=0\big)
= \frac{1}{N_1 N_0} \sum_{i:\,y_i=1}\ \sum_{j:\,y_j=0}
\Big[ \mathbb{1}(\hat{p}_i > \hat{p}_j) + \tfrac{1}{2}\,\mathbb{1}(\hat{p}_i = \hat{p}_j) \Big]
\end{equation}
Es insensible al umbral y a la escala monot\'onica de $\hat{p}$; por eso el
desbalance 8.1\% no obliga a remuestrear.

\subsection{Gradient boosting (formulaci\'on com\'un)}
Los tres modelos base son m\'aquinas de boosting aditivo sobre \'arboles:
\begin{equation}
\hat{p}_i = \sigma\!\left( \eta_i \right), \qquad
\eta_i^{(T)} = \eta_i^{(0)} + \sum_{t=1}^{T} \nu\, f_t(x_i), \qquad
\sigma(z) = \frac{1}{1+e^{-z}}
\end{equation}
con tasa de aprendizaje $\nu = 0.03$. En la iteraci\'on $t$, $f_t$ se ajusta
minimizando la p\'erdida log\'istica:
\begin{equation}
\mathcal{L} = -\frac{1}{N}\sum_{i=1}^{N}
\Big[ y_i \log \hat{p}_i + (1-y_i)\log(1-\hat{p}_i) \Big].
\end{equation}

\subsection{XGBoost: objetivo de segundo orden}
Aproximaci\'on de Taylor de orden 2 del objetivo regularizado:
\begin{equation}
\mathcal{L}^{(t)} \simeq \sum_{i=1}^{N}\Big[ g_i f_t(x_i) + \tfrac{1}{2} h_i f_t^2(x_i) \Big]
+ \gamma T_{\text{hojas}} + \tfrac{1}{2}\lambda \sum_{j} w_j^2,
\quad g_i = \hat{p}_i - y_i,\quad h_i = \hat{p}_i(1-\hat{p}_i)
\end{equation}
y el peso \'optimo de la hoja $j$ y la ganancia de un split:
\begin{equation}
w_j^* = -\frac{G_j}{H_j + \lambda}, \qquad
\mathrm{Gain} = \frac{1}{2}\left[ \frac{G_L^2}{H_L+\lambda} + \frac{G_R^2}{H_R+\lambda}
- \frac{(G_L+G_R)^2}{H_L+H_R+\lambda} \right] - \gamma .
\end{equation}

\subsection{LightGBM: histogramas y GOSS/EFB}
Mismo objetivo, pero discretiza las features en histogramas de $b$ bins
(reduce coste de split de $O(N)$ a $O(b)$ por feature) y crece hojas
\emph{leaf-wise} con l\'imites \texttt{num\_leaves}$=31$ y
\texttt{min\_child\_samples}$=50$. La regularizaci\'on expl\'icita del
baseline es $\texttt{lambda\_l2} = 1.0$ sobre los pesos de hoja, m\'as
submuestreo estoc\'astico: \texttt{feature\_fraction}$=0.8$ por \'arbol y
\texttt{bagging\_fraction}$=0.8$ por iteraci\'on.

\subsection{CatBoost: target statistics ordenadas}
Para la categ\'orica $k$ de la fila $i$, el encoding se calcula solo con filas
\emph{anteriores} en una permutaci\'on aleatoria (anti-leakage):
\begin{equation}
\mathrm{TS}(x_i^{k}) = \frac{\sum_{j \prec i} \mathbb{1}(x_j^{k} = x_i^{k})\, y_j + a\, \bar{y}}
{\sum_{j \prec i} \mathbb{1}(x_j^{k} = x_i^{k}) + a}
\end{equation}
con prior global $\bar{y}$ y smoothing $a$. Su an\'alogo manual en el proyecto
(target encoding OOF con smoothing) usa la media de los otros folds como
denominador de evidencia.

\subsection{Target encoding OOF con smoothing (features manuales)}
Para \texttt{ORGANIZATION\_TYPE} / \texttt{OCCUPATION\_TYPE}, fold $k$ excluido:
\begin{equation}
\mathrm{TE}(c) = \frac{n_c\, \bar{y}_c + m\, \bar{y}}{n_c + m}
\end{equation}
donde $\bar{y}_c$ y $n_c$ son media y conteo de la categor\'ia $c$ \textbf{en
los folds de entrenamiento}, y $m$ controla el encogimiento al prior. La fila
nunca contribuye a su propio encoding (sin leakage).

\subsection{Validaci\'on cruzada: estimador CV y RSKF}
StratifiedKFold preserva la proporci\'on de positivos en cada fold
($\pi_k \approx 0.081$). El estimador de control diario:
\begin{equation}
\widehat{\mathrm{AUC}}_{\mathrm{OOF}} = \mathrm{AUC}\big(\{(y_i, \hat{p}_i^{\,\mathrm{OOF}})\}_{i=1}^{N}\big)
\end{equation}
Para la decisi\'on final, Repeated StratifiedKFold $3 \times 5$ (semillas
42/43/44) da $R=15$ estimaciones y la decisi\'on usa
\begin{equation}
\bar{A} = \frac{1}{R}\sum_{r=1}^{R} A_r, \qquad
s_A = \sqrt{\frac{1}{R-1}\sum_{r=1}^{R} (A_r - \bar{A})^2} .
\end{equation}
El candidato gana solo si $\bar{A}_{\text{cand}} > \bar{A}_{\text{ref}} + \delta$
con $\delta \in \{0.003\ \text{features},\ 0.002\ \text{tuning},\ 0.001\
\text{ensemble}\}$.

\subsection{Blend ponderado (nivel 1)}
Con predicciones OOF de $M$ modelos, el blend resuelve:
\begin{equation}
w^* = \arg\max_{w \in \Delta} \mathrm{AUC}\!\left( \left\{ y_i,\ \textstyle\sum_{m=1}^{M} w_m\, \hat{p}_{i,m}^{\mathrm{OOF}} \right\} \right),
\quad \Delta = \{ w_m \ge 0,\ \textstyle\sum_m w_m = 1 \}
\end{equation}
explorado en rejilla con paso $0.1$. Como AUC es invariante monot\'onico, no se
impone calibraci\'on previa.

\subsection{Stacking (nivel 2, Sprint 4)}
Meta-modelo log\'istico sobre las OOF apiladas:
\begin{equation}
\hat{p}_i^{\text{stack}} = \sigma\!\left( \beta_0 + \sum_{m=1}^{M} \beta_m\, \mathrm{logit}\, \hat{p}_{i,m}^{\mathrm{OOF}} \right),
\qquad \beta^* = \arg\min_\beta \sum_i \ell\big(y_i, \hat{p}_i^{\text{stack}}\big) + \lambda_\beta \|\beta\|_2^2
\end{equation}
entrenado con la misma validaci\'on anidada para no inflar el estimador.

\section{Metas}
\begin{itemize}
  \item \textbf{M3 (6 nov) --- red de seguridad:} pipeline end-to-end
        reproducible (\texttt{python -m src.train}) con CV $\geq 0.7681$
        (baseline $+0.01$), \texttt{refs.md} completo.
  \item \textbf{M5 (18 nov):} submission final elegida por RSKF, dos d\'ias de
        colch\'on antes del cierre oficial.
\end{itemize}

\end{document}
